% Standalone problem document, generated from the adjacent README.md.
% Compile with XeLaTeX. Mathematical content is maintained in Markdown.
\documentclass[11pt,a4paper]{article}
\usepackage[fontsize=10.5pt]{scrextend}
\usepackage[margin=22mm,headheight=15pt,headsep=9mm,footskip=11mm]{geometry}
\usepackage{fontspec}
\setmainfont{texgyrepagella}[Extension=.otf,UprightFont=*-regular,BoldFont=*-bold,ItalicFont=*-italic,BoldItalicFont=*-bolditalic]
\setsansfont{texgyreheros}[Extension=.otf,UprightFont=*-regular,BoldFont=*-bold,ItalicFont=*-italic,BoldItalicFont=*-bolditalic]
\setmonofont{texgyrecursor}[Extension=.otf,UprightFont=*-regular,BoldFont=*-bold,ItalicFont=*-italic,BoldItalicFont=*-bolditalic,Scale=MatchLowercase]
\usepackage{amsmath,amssymb}
\usepackage{unicode-math}
\setmathfont{texgyrepagella-math.otf}
\usepackage{xcolor,microtype,enumitem,needspace,fancyhdr}
\usepackage{bookmark}
\definecolor{ink}{HTML}{17344B}
\definecolor{accent}{HTML}{197D80}
\definecolor{muted}{HTML}{596773}
\hypersetup{colorlinks=true,linkcolor=accent,urlcolor=accent,pdftitle={MF-19 - Decidability
of zero hitting by rational matrix
powers},pdfauthor={Open Problems in Numerical Linear Algebra}}
\urlstyle{same}
\setlength{\parindent}{0pt}
\setlength{\parskip}{5pt plus 1pt minus 1pt}
\linespread{1.04}
\setlength{\emergencystretch}{3em}
\widowpenalty=10000
\clubpenalty=10000
\displaywidowpenalty=10000
\allowdisplaybreaks[1]
\setlist{nosep,leftmargin=1.5em}
\providecommand{\tightlist}{\setlength{\itemsep}{0pt}\setlength{\parskip}{0pt}}
\providecommand{\pandocbounded}[1]{#1}
\pagestyle{fancy}
\fancyhf{}
\fancyhead[L]{\sffamily\footnotesize\color{muted}OPEN PROBLEMS IN NUMERICAL LINEAR ALGEBRA}
\fancyhead[R]{\sffamily\bfseries\footnotesize\color{accent}MF-19}
\fancyfoot[L]{\parbox[b]{0.9\textwidth}{\sffamily\fontsize{8}{10}\selectfont\color{muted}Editorial ratings; a bounded search cannot certify that a problem remains unsolved.\\Literature check: 2026-09-10}}
\fancyfoot[R]{\sffamily\footnotesize\color{muted}\thepage}
\renewcommand{\headrulewidth}{0.3pt}
\renewcommand{\headrule}{\hbox to\headwidth{\color{accent}\leaders\hrule height \headrulewidth\hfill}}
\makeatletter
\renewcommand{\section}{\Needspace{5\baselineskip}\@startsection{section}{1}{0pt}{12pt plus 2pt minus 2pt}{4pt}{\sffamily\bfseries\large\color{ink}}}
\renewcommand{\subsection}{\Needspace{4\baselineskip}\@startsection{subsection}{2}{0pt}{9pt plus 1pt minus 1pt}{3pt}{\sffamily\bfseries\normalsize\color{ink}}}
\makeatother
\setcounter{secnumdepth}{0}
\begin{document}
{\sffamily\small\color{muted}Matrix functions and stability\par}
\vspace{3pt}
{\raggedright\hyphenpenalty=10000\exhyphenpenalty=10000\sffamily\bfseries\fontsize{21}{24}\selectfont\color{ink}Decidability
of zero hitting by rational matrix powers\par}
\vspace{7pt}
{\sffamily\small\textbf{Difficulty:} extreme\quad\textbf{Importance:} broadly
interesting\par}
{\sffamily\small\bfseries\color{accent}Status: Partially resolved\par}
\vspace{4pt}
{\color{accent}\hrule height 0.8pt}
\vspace{6pt}
\textbf{Rating rationale:} Unconditional decidability in unrestricted
dimension is an extreme longstanding barrier; reachability, recurrences
and exact linear dynamics give broad significance.

\subsection{Problem statement}\label{problem-statement}

Does there exist a single algorithm which, given a positive integer
\(d\), a matrix \(A\in\mathbb Q^{d\times d}\), and vectors
\(u,v\in\mathbb Q^d\), always halts and correctly decides whether

\[
u^{T}A^kv=0\qquad\text{for some integer }k\ge0?
\]

All rational inputs are represented exactly, for example by binary
integers for numerators and nonzero denominators. There is no a priori
bound on \(d\), on the entries, or on the hitting time \(k\). The output
sought is a yes/no answer; numerical near-zero detection does not
suffice.

This is the Skolem problem in its matrix-power form. Equivalently, one
is given rational coefficients and initial values for a finite-order
scalar linear recurrence and asks whether one of its terms vanishes.
These equivalent formulations form one catalog entry.

The matrix formulation asks whether a discrete linear system reaches a
specified hyperplane. It is directly relevant to exact observability,
reachability, and verification of linear dynamics, and exposes a
fundamental limit of computations with powers of matrices.

\subsection{References}\label{references}

\begin{enumerate}
\def\labelenumi{\arabic{enumi}.}
\tightlist
\item
  P. Bacik, T. Karimov, F. Luca, J. Nieuwveld, J. Ouaknine, D. Purser
  and J. Worrell, \emph{A survey of the Skolem and Positivity Problems
  for linear recurrence sequences}, submitted (2026), §1, Problem 2, and
  the discussion of linear dynamical systems; §5.
  \href{https://people.mpi-sws.org/~joel/publications/skolem_and_positivity_survey26.pdf}{Author-hosted
  PDF}.
\item
  F. Luca, J. Ouaknine and J. Worrell, \emph{Conjectural Decidability of
  the Skolem Problem}, arXiv:2607.15510v1 (2026), §1.
  \href{https://arxiv.org/pdf/2607.15510}{Full preprint}.
\end{enumerate}

Status check (2026-09-10): the 2026 survey still poses the general
problem. The July 16 preprint proves conditional decidability under a
strengthened prime-gap conjecture and an unconditional density-one
universal Skolem-set result; it does not supply an unconditional
algorithm on all inputs. Its current arXiv record is v1. Searches for
``Skolem Problem decidability solution 2026'', including September
developments, found special-family and positive-characteristic results
but no general rational-input resolution. This is a dated, bounded
status check.

\subsection{Audit --- 2026-09-10}\label{audit-2026-09-10}

Independently rechecked
\href{https://arxiv.org/pdf/2607.15510}{Luca--Ouaknine--Worrell, §1},
its current v1 record, and the
\href{https://people.mpi-sws.org/~joel/publications/skolem_and_positivity_survey26abs.html}{2026
survey's author abstract}. The order-at-most-four decidability result
gives a substantive exact subclass (in particular, \(d\le4\) here),
justifying Partially resolved; conditional decidability and density-one
sets do not settle unrestricted inputs. Later Skolem-decidability
searches found no unconditional general resolution. Both ratings are
retained.
\end{document}
